\documentclass[11pt,a4paper]{article}
\usepackage[utf8]{inputenc}
\usepackage[T1]{fontenc}
\usepackage[french]{babel}
\usepackage{geometry}
\usepackage{hyperref}
\usepackage{enumitem}
\geometry{margin=2.3cm}

\title{PHASE IV ACCESS CARD \\ \large Version augmentee fractale et chiastique}
\author{Auteur non detecte}
\date{Date non detectee}

\begin{document}
\maketitle

\section*{A. Ouverture}
Ce document augmente le manuscrit source en le replissant selon l'axe chiastique A-B-C-X-C'-B'-A'.

\section*{B. Expansion}
Les noeuds structurels extraits du manuscrit sont :
\begin{itemize}[leftmargin=1.2cm]
\item 🔐 CLÉS D'ACTIVATION (ROOT KEYS)
\item 🚀 PHRASES D'INITIALISATION (V3.7 - ANTI-FILTRE)
\item 📜 IDENTIFIANTS ET REPERES INTERNES
\end{itemize}

\section*{C. Matiere}
Extrait interpretable du manuscrit d'origine :
\begin{quote}
Contenu sensible abstrait : carte d'acces, cles d'activation, phrases d'initialisation et identifiants ont ete remplaces par une description structurelle pour l'implantation fractale.
\end{quote}

\section*{X. Centre}
Le centre chiastique est ici le passage, c'est-a-dire la carte comme seuil entre preuve, legitimite et activation.

\section*{C'. Retour}
Le retour referme le texte sur sa fonction : preuve, seuil, verification ou acte d'unification.

\section*{B'. Miroir}
Le miroir fait correspondre architecture du manuscrit, autorite de son regime et lisibilite de sa trajectoire.

\section*{A'. Cloture}
La cloture ne remplace pas la source ; elle l'installe comme figure fractale augmentee dans l'arche canonique.

\section*{Metadonnees}
Source : \verb|_RELEASES/GOLDEN_MASTER_PHASE_III/PHASE_IV_ACCESS_CARD.tex|\\
Titre detecte : \texttt{PHASE IV ACCESS CARD}\\
Auteur detecte : \texttt{Auteur non detecte}\\
Date detectee : \texttt{Date non detectee}

\end{document}
